\section{Compactness of continuous homomorphisms}

The compactness step in \cite{Camina1997}, Theorem~5, can be expressed
entirely in finite quotients. The argument below replaces the paper's
evaluation of infinite pro-$p$ words and its appeal to automatic
continuity by a diagonal subsequence of continuous quotient maps.
The required depth control on the finite embeddings remains an
explicit hypothesis of the final lemma in this section.

\begin{lemma}[Finite targets and finitely many generators]
\label{lem:finite-continuous-hom-set}
\uses{def:topologically-finitely-generated}
\source{camina1997:page-0011}
\notready
Let $P$ be a topological group with a finite tuple
$(s_1,\ldots,s_r)$ generating a dense subgroup. If $A$ is a finite
discrete group, there are at most $|A|^r$ continuous homomorphisms
$P\to A$.
\end{lemma}

\begin{proof}
Two homomorphisms with the same values on the tuple agree on the
subgroup it generates, by preservation of products and inverses.
If both are continuous, their equalizer is closed: it is the
inverse image of the diagonal under their product map, and the
diagonal in $A\times A$ is closed because $A$ is Hausdorff.
The equalizer contains a dense subgroup and hence all of $P$.
Thus evaluation at the tuple injects the continuous homomorphism
set into $A^r$, which has $|A|^r$ elements.
\end{proof}

\begin{lemma}[Countability of the finite quotient basis]
\label{lem:finitely-generated-pro-p-countable-base}
\uses{def:pro-p-group,def:countably-based,lem:finite-continuous-hom-set}
\notready
Every topologically finitely generated pro-$p$ group is countably based.
\end{lemma}

\begin{proof}
For each integer $d>0$, there are only finitely many group laws
on a fixed set of size $d$. By the preceding lemma there are only
finitely many continuous homomorphisms from $P$ to each such group.
Consequently the kernels of continuous homomorphisms to finite
discrete groups form a countable family, by taking the union over $d$.
Every open normal subgroup of finite index is such a kernel:
the quotient map is continuous because its fibres are open cosets,
and one may identify the quotient with a group on a finite labelled set.
The pro-$p$ definition supplies an identity neighbourhood basis
of open normal subgroups of finite index. That basis is contained
in the countable family just described, whose members are themselves
open identity neighbourhoods. Thus this family is a countable basis.
\end{proof}

\begin{lemma}[A diagonal limit of continuous Nottingham homomorphisms]
\label{lem:nottingham-hom-diagonal-limit}
\uses{lem:finite-continuous-hom-set,lem:nottingham-quotient-card,lem:nottingham-inverse-limit}
\source{camina1997:page-0011,camina1997:page-0012}
\notready
Let $P$ be a topologically finitely generated topological group and
let $f_i:P\to J_p$, for $i\in\mathbf N$, be continuous homomorphisms.
There are a strictly increasing sequence $(i_j)$ and a continuous
homomorphism $f:P\to J_p$ such that, for every $n\ge0$, the maps
\[
 P\xrightarrow{f_{i_j}}J_p\longrightarrow J_p/J_{p,n}
\]
are equal to the quotient of $f$ for all sufficiently large $j$.
The threshold may be chosen independently of the element of $P$.
\end{lemma}

\begin{proof}
Write $Q_n=J_p/J_{p,n}$ and let $\pi_n:J_p\to Q_n$ be the
continuous quotient map. The set of continuous homomorphisms
$P\to Q_n$ is finite by the first lemma and finiteness of $Q_n$.
Starting with all indices, apply the infinite pigeonhole principle
at level $0$ to choose an infinite set $S_0$ on which
$\pi_0f_i$ is constant, with value $h_0$. Inductively, partition
the infinite set $S_n$ by the finitely many values of
$\pi_{n+1}f_i$ and choose an infinite part $S_{n+1}\subseteq S_n$,
with constant value $h_{n+1}$.
Choose $i_0\in S_0$ and, recursively, choose
$i_{j+1}\in S_{j+1}$ larger than $i_j$. Such a choice exists
because an infinite subset of $\mathbf N$ is unbounded.
For $j\ge n$, nesting gives $i_j\in S_n$, so
$\pi_nf_{i_j}=h_n$ as maps on all of $P$.

The maps $h_n$ are compatible. Indeed, for $m\ge n$, choose
$j\ge m$. Both $h_m$ and $h_n$ are quotients of the same map
$f_{i_j}$, so $\pi_{m,n}h_m=h_n$, where $\pi_{m,n}:Q_m\to Q_n$
is the natural quotient homomorphism.
For each $x\in P$, the tuple $(h_n(x))_n$ therefore lies in
the congruence inverse limit. It depends homomorphically on $x$,
since every coordinate does. It also depends continuously on $x$:
all coordinates are continuous, and product continuity followed
by restriction to the inverse-limit subspace gives continuity.
Compose this map with the inverse of the topological group
isomorphism $J_p\cong\varprojlim Q_n$. This gives the required
continuous homomorphism $f$. By construction $\pi_nf=h_n$, and
the equality for all $j\ge n$ gives the claimed uniform threshold.
\end{proof}

\begin{lemma}[Depth control prevents collapse in the limit]
\label{lem:nottingham-depth-controlled-limit-embedding}
\uses{lem:nottingham-hom-diagonal-limit,lem:compact-injection-closed-embedding}
\source{camina1997:page-0012}
\notready
Let $P$ be a compact, topologically finitely generated topological
group. Suppose that $(f_i:P\to J_p)_{i\in\mathbf N}$ is a sequence
of continuous homomorphisms such that for every $x\ne1$ there
exist $n,N\in\mathbf N$ with
\[
 f_i(x)\notin J_{p,n}\qquad\text{for every }i\ge N.
\]
Then $P$ admits a closed group embedding in $J_p$.
\end{lemma}

\begin{proof}
Take $f$ and the strictly increasing subsequence $(i_j)$ from
the preceding lemma. Suppose $x\ne1$, and choose $n,N$ as in
the hypothesis. For all sufficiently large $j$, we have
$i_j\ge N$ and $\pi_nf_{i_j}=\pi_nf$. The first assertion holds
because a strictly increasing sequence of natural numbers is
unbounded; the second is the diagonal-limit conclusion.
The hypothesis gives $\pi_nf_{i_j}(x)\ne1$, and hence
$\pi_nf(x)\ne1$. Thus $f(x)\ne1$. Its kernel is trivial, so
$f$ is injective. It is continuous by its construction, $P$ is
compact by hypothesis, and $J_p$ is Hausdorff. The compact-source
lemma therefore makes $f$ a homeomorphism onto a closed subgroup.
\end{proof}

In Camina's tower, each $f_i$ is a finite-quotient map followed by
an embedding of that finite quotient in $J_p$. Her depth comparison
makes the depths of $f_i(x)$ eventually constant for each $x\ne1$;
this implies the avoidance hypothesis above by choosing a deeper
congruence subgroup. Constructing those finite embeddings with
simultaneous depth control is the separate ramification step.
